\section{Conclusion and limitations}
\label{sec:discussion}

We described compositionality as a spectrum of capabilities with two forms. Weak compositionality recombines known parts. Strong compositionality keeps every weak capability and adds decomposition, the return of each part in its role from a whole. Decomposition becomes a useful notion only under a simple class of procedures on held-out parts, and recoverability scores it under both conditions. With these terms as a shared language, we placed fifteen accounts of compositionality on one spectrum.

The evidence supports the hypothesis that behavior is a subset of representation. Language models that answered by meaning carried the role in the state at the question, and patches of that state changed their answers. Glyph networks, recursive composers and image encoders kept parts in their states that their own procedures lost. Procedures that read those states reached the strong capabilities the systems lacked. OLMo-2 7B answered beyond our readouts, which marks a limit of the declared class, whereas models that answered by word order failed in both places. Pretrained models built recoverability beyond their inputs, but small trained networks inherited it from theirs.

Our results suggest five requirements for a strongly compositional network. Its state must have room for every part (Propositions~\ref{prop:ceiling} and~\ref{prop:ey}). Its procedure for reading parts must learn from many distinct parts (Proposition~\ref{prop:span}). Each part should enter every role through one code, so that one composite fixes its reversal (Proposition~\ref{prop:systematic}). The system must use those role codes. Finally, the relation must stay invertible on its own outputs, and the system's procedure must follow it (Proposition~\ref{prop:depth}). We therefore argue that a claim of compositionality should name its form and its rung. It should also declare the readout class, part codes and held-out parts behind it.

\paragraph{Limitations and next steps.} The inclusion of behavior in representation is a logical result for a readout class that contains the system's procedure. The evidence that the inclusion is strict comes from the systems we tested. Our evidence for design comes from repairs and from a recovery objective that also adds data, so separating the objective from the data is the next test. Our trained networks are small, with rendered domains so that we could hold parts out, and their recoverability comes largely from their inputs. In pretrained models we hold parts out only from the readout, apart from invented words. Richer methods, such as a discovered tensor-product decomposition \citep{mccoy2019}, may find structure that our readouts miss. A human study could ask whether people show the same gap, by testing whether they recover the parts of composites built from newly learned parts.

